\documentclass[11pt]{article}

\usepackage[T1]{fontenc}
\usepackage{lmodern}
\usepackage{amsmath,amssymb,amsthm,mathtools}
\usepackage[margin=1in]{geometry}
\usepackage{microtype}
\usepackage[hidelinks]{hyperref}

\newtheorem{lemma}{Lemma}
\newtheorem{proposition}[lemma]{Proposition}
\newtheorem{corollary}[lemma]{Corollary}
\theoremstyle{remark}
\newtheorem{remark}[lemma]{Remark}

\newcommand{\N}{\mathbb N}
\newcommand{\R}{\mathbb R}
\newcommand{\ind}{\mathbf 1}
\newcommand{\strict}[1]{\left[#1\right]_{<}}

\title{A Global Structural Reduction for the Spherical-Shell Proxy}
\author{}
\date{Global analytic structural module}

\begin{document}
\maketitle

\begin{abstract}
This note gives the global structural component of the purely analytic
middle-ratio argument for the spherical-shell proxy.  The weighted proxy is
first written as an exact layer-cake sum.  Its deficit from the action
integral is then decomposed into a radial shifted-quadrature deficit and an
angular rounding defect.  A Stieltjes calculation retains two nonnegative
remainders that are discarded by the optical estimate, and one of those
remainders admits a smooth explicit lower bound.  The result is the direct
lower bound \eqref{eq:direct-structural-gap}.  All structural identities
hold for every $0<\rho<1$ and $K>0$; in particular they apply on the
upward-closed
middle-ratio region
\(
 \rho_c\leq\rho\leq39/50,
 \ K\geq k_{11}(\rho)
\).
\end{abstract}

\section{Action and the global target}

For $a>0$ and $z\geq0$, put
\begin{equation}
 G_a(z)=
 \begin{cases}
 \displaystyle \frac1\pi\left(\sqrt{a^2-z^2}
       -z\arccos\frac za\right),&0\leq z<a,\\[4pt]
 0,&z\geq a.
 \end{cases}
 \label{eq:G-def}
\end{equation}
Fix $0<\rho<1$ and $K>0$, and define
\begin{equation}
 A(z)=A_{\rho,K}(z):=G_K(z)-G_{\rho K}(z),
 \qquad
 T:=A(0)=\frac{(1-\rho)K}{\pi}.
 \label{eq:A-def}
\end{equation}
The strict positive-integer count is
\[
 \strict{x}:=\#\{n\in\N:n<x\}.
\]
The spherical-shell proxy and its continuous payment are
\begin{align}
 P(\rho,K)
 &:=\sum_{\substack{\nu=\ell+1/2<K\\ \ell\in\mathbb N_0}}
       2\nu\,\strict{A(\nu)+\tfrac14},
 \label{eq:proxy-def}\\
 W(\rho,K)
 &:=\int_0^K2zA(z)\,dz
   =\frac{2(1-\rho^3)K^3}{9\pi}.
 \label{eq:W-def}
\end{align}
The upward-closed middle-ratio region used below is
\begin{equation}
 \mathcal R_{\mathrm c}:=
 \left\{(\rho,K):
 \rho_c\leq\rho\leq\frac{39}{50},\quad
 K\geq k_{11}(\rho)\right\},
 \qquad
 \rho_c:=\frac1{1+2\pi},
 \qquad
 k_{11}(\rho):=
 \sqrt{\frac{\pi^2}{(1-\rho)^2}+132}.
 \label{eq:compact-residual}
\end{equation}
The middle-ratio target is $P<W$ on $\mathcal R_{\mathrm c}$.

For $0<z<K$, direct differentiation gives
\begin{equation}
 -A'(z)=\frac1\pi
 \begin{cases}
 \displaystyle
 \arccos\frac zK-\arccos\frac z{\rho K},&0<z<\rho K,\\[6pt]
 \displaystyle \arccos\frac zK,&\rho K\leq z<K.
 \end{cases}
 \label{eq:A-prime}
\end{equation}
In particular, $A$ is strictly decreasing from $T$ to $0$.  Let
\[
 R:[0,T]\longrightarrow[0,K]
\]
be its decreasing inverse, and set
\begin{equation}
 F(t):=R(t)^2,
 \qquad
 g(t):=-F'(t),
 \qquad
 \tau:=A(\rho K)=K G_1(\rho).
 \label{eq:inverse-data}
\end{equation}
The derivative in \eqref{eq:A-prime} increases in magnitude on
$(0,\rho K)$ and decreases in magnitude on $(\rho K,K)$.  The inverse
curvature has the related, but slightly stronger, property
\begin{equation}
 g\text{ decreases on }(0,\tau),
 \qquad
 g\text{ increases on }(\tau,T).
 \label{eq:g-shape}
\end{equation}
Indeed, on the outer branch, if $z=R(t)$ then
\[
 g(t)=\frac{2\pi z}{\arccos(z/K)};
\]
the quotient increases with $z$, while $z$ decreases with $t$.  On the
inner branch, put
\[
 \delta(z):=\arccos\frac zK-\arccos\frac z{\rho K}.
\]
Differentiation with respect to the intermediate ratio gives
\[
 \frac{\delta(z)}z
 =\int_\rho^1\frac{ds}{s\sqrt{s^2K^2-z^2}}.
\]
This integral increases with $z$.  Thus $z/\delta(z)$ decreases with
$z$, and $g(t)=2\pi z/\delta(z)$ increases as $t$ increases and $z$
decreases.
The endpoint and interface data are
\begin{equation}
 F(\tau)=\rho^2K^2,
 \qquad
 g(\tau)=\frac{2\pi\rho K}{\arccos\rho},
 \qquad
 g(T)=\frac{2\pi\rho K}{1-\rho}.
 \label{eq:g-data}
\end{equation}
Here and below $g(T)$ is the one-sided limit.  Near $t=0$, one has
$g(t)=O(t^{-1/3})$, so all improper Stieltjes boundary terms used below
vanish.

\section{Exact layer cake and angular rounding}

Let
\begin{equation}
 x_n:=n-\frac14,
 \qquad
 M(r):=\#\left\{\ell\in\mathbb N_0:\ell+\frac12<r\right\}
       =\max\left\{0,\left\lceil r-\frac12\right\rceil\right\}.
 \label{eq:shifted-grids}
\end{equation}

\begin{lemma}[Exact layer cake]
\label{lem:layer-cake}
For every $0<\rho<1$ and $K>0$,
\begin{equation}
 P(\rho,K)=\sum_{\substack{n\geq1\\x_n<T}}M(R(x_n))^2,
 \qquad
 W(\rho,K)=\int_0^T F(t)\,dt.
 \label{eq:layer-cake}
\end{equation}
Both identities retain the strict convention at every radial and angular
wall.
\end{lemma}

\begin{proof}
For a fixed half-integer $\nu$,
\[
 \strict{A(\nu)+\frac14}
 =\#\{n\geq1:x_n<A(\nu)\}.
\]
Since $A$ and $R$ are decreasing, $x_n<A(\nu)$ is equivalent to
$\nu<R(x_n)$.  Interchanging the two finite sums and using
\[
 \sum_{\ell=0}^{m-1}(2\ell+1)=m^2
\]
proves the first identity.  At $R(x_n)=m+1/2$, the strict angular count is
$m$, exactly as specified by $M$.

For the second identity, the area formula for a decreasing inverse gives
\[
 \int_0^T R(t)^2\,dt
 =\int_0^K2zA(z)\,dz.
\]
The last integral is \eqref{eq:W-def}.
\end{proof}

For every active $n$, abbreviate
\begin{equation}
 r_n:=R(x_n),
 \qquad m_n:=M(r_n),
 \qquad
 \alpha_n:=r_n-m_n+\frac12.
 \label{eq:angular-phase}
\end{equation}
The strict convention implies
\begin{equation}
 m_n-\frac12<r_n\leq m_n+\frac12,
 \qquad 0<\alpha_n\leq1.
 \label{eq:angular-cell}
\end{equation}
This includes $m_n=0$.  Define the angular rounding defect
\begin{equation}
 \mathcal E_{\mathrm{ang}}(\rho,K)
 :=\sum_{x_n<T}(m_n^2-r_n^2).
 \label{eq:E-ang-def}
\end{equation}
The defect of a single radial layer is exactly
\begin{equation}
 m_n^2-r_n^2
 =(1-2\alpha_n)r_n
  +\left(\alpha_n-\frac12\right)^2.
 \label{eq:E-ang-cell}
\end{equation}
Unlike the one-sided estimate
$m_n^2-r_n^2<r_n+1/4$, formula
\eqref{eq:E-ang-cell} retains the cancellation when
$\alpha_n>1/2$.

Combining \eqref{eq:layer-cake} and \eqref{eq:E-ang-def} gives the first
exact two-term decomposition:
\begin{equation}
 W-P=\mathcal D_{\mathrm{rad}}-\mathcal E_{\mathrm{ang}},
 \qquad
 \mathcal D_{\mathrm{rad}}
 :=\int_0^TF(t)\,dt-\sum_{x_n<T}F(x_n).
 \label{eq:two-term-basic}
\end{equation}

\section{The radial deficit with retained positive remainders}

Define the shifted counting function and its two primitives by
\begin{equation}
 C(t):=\#\{n\geq1:x_n<t\},
 \quad
 w(t):=t-C(t),
 \quad
 \mathfrak W(t):=\int_0^tw(s)\,ds,
 \quad
 \Psi(t):=\mathfrak W(t)-\frac t4.
 \label{eq:sawtooth-def}
\end{equation}
On $0\leq t\leq3/4$,
\begin{equation}
 \Psi(t)=\frac12\left(t-\frac14\right)^2-\frac1{32}.
 \label{eq:first-cell}
\end{equation}
For $t=x_n+u$, $0\leq u\leq1$,
\begin{equation}
 \Psi(t)=\frac12\left(u-\frac12\right)^2-\frac1{32}.
 \label{eq:periodic-cell}
\end{equation}
Consequently
\begin{equation}
 \mathfrak W(t)\geq0,
 \qquad
 -\frac1{32}\leq\Psi(t)\leq\frac3{32}
 \quad(t\geq0).
 \label{eq:sawtooth-bounds}
\end{equation}

\begin{proposition}[Exact retained-remainder identity]
\label{prop:exact-deficit}
For every $0<\rho<1$ and $K>0$, put
\begin{align}
 \mathcal R_{\mathrm{dec}}
 &:=-\int_{(0,\tau)}\mathfrak W(t)\,dg(t),
 \label{eq:R-dec}\\
 \mathcal R_{\mathrm{osc}}
 &:={}
 \left(\Psi(T)+\frac1{32}\right)g(T)
 +\int_{(\tau,T)}
       \left(\frac3{32}-\Psi(t)\right)\,dg(t).
 \label{eq:R-osc}
\end{align}
Then both remainders are nonnegative, and
\begin{equation}
 \begin{aligned}
 \mathcal D_{\mathrm{rad}}
 ={}&\frac{F(\tau)}4-\frac{g(T)}8
 +g(\tau)\left(\frac\tau4+\frac3{32}\right)\\
 &+\mathcal R_{\mathrm{dec}}+\mathcal R_{\mathrm{osc}}.
 \end{aligned}
 \label{eq:exact-radial-deficit}
\end{equation}
Therefore
\begin{equation}
 W-P=
 \mathcal B(\rho,K)
 +\mathcal R_{\mathrm{dec}}+\mathcal R_{\mathrm{osc}}
 -\mathcal E_{\mathrm{ang}},
 \label{eq:master-identity}
\end{equation}
where
\begin{equation}
 \mathcal B(\rho,K):=
 \frac{\rho^2K^2}{4}
 -\frac{\pi\rho K}{4(1-\rho)}
 +\frac{2\pi\rho K}{\arccos\rho}
  \left(\frac{KG_1(\rho)}4+\frac3{32}\right).
 \label{eq:B-def}
\end{equation}
\end{proposition}

\begin{proof}
Distributionally, $dw=dt-dC$.  Stieltjes integration by parts, using
$F(T)=0$ and $w(0)=0$, gives
\begin{equation}
 \mathcal D_{\mathrm{rad}}=\int_0^T w(t)g(t)\,dt.
 \label{eq:radial-stieltjes}
\end{equation}
The improper boundary term at $0$ vanishes because
$\mathfrak W(t)=t^2/2$ there and $g(t)=O(t^{-1/3})$.

On $(0,\tau)$, integrate with $d\mathfrak W=w\,dt$:
\begin{equation}
 \int_0^\tau wg
 =\mathfrak W(\tau)g(\tau)
  -\int_{(0,\tau)}\mathfrak W\,dg.
 \label{eq:decreasing-part}
\end{equation}
On $(\tau,T)$, use $w=1/4+\Psi'$ and
$\int_\tau^Tg=F(\tau)$:
\begin{equation}
 \int_\tau^T wg
 =\frac{F(\tau)}4
  +\Psi(T)g(T)-\Psi(\tau)g(\tau)
  -\int_{(\tau,T)}\Psi\,dg.
 \label{eq:increasing-part}
\end{equation}
Since $\mathfrak W(\tau)=\tau/4+\Psi(\tau)$, adding
\eqref{eq:decreasing-part} and \eqref{eq:increasing-part} cancels the
uncontrolled interface value $\Psi(\tau)$.

Now add and subtract the extremal sawtooth values in
\eqref{eq:sawtooth-bounds}.  The exact identity
\begin{align*}
 &\Psi(T)g(T)-\int_{(\tau,T)}\Psi\,dg\\
 &\quad={}-\frac{g(T)}8+\frac{3g(\tau)}{32}
 +\left(\Psi(T)+\frac1{32}\right)g(T)
 +\int_{(\tau,T)}
     \left(\frac3{32}-\Psi\right)dg
\end{align*}
gives \eqref{eq:exact-radial-deficit}.  By
\eqref{eq:g-shape}, $-dg$ is a positive measure on $(0,\tau)$ and
$dg$ is a positive measure on $(\tau,T)$.  Equations
\eqref{eq:sawtooth-bounds}, \eqref{eq:R-dec}, and
\eqref{eq:R-osc} therefore prove nonnegativity of both remainders.
Finally substitute \eqref{eq:g-data} and $\tau=KG_1(\rho)$, and then use
\eqref{eq:two-term-basic}.
\end{proof}

\section{A smooth lower bound for the decreasing-branch remainder}

The term $\mathcal R_{\mathrm{dec}}$ is quantitatively important on the
middle-ratio region; discarding it loses too much.  It has the following
elementary smooth minorant.

\begin{lemma}[A global primitive minorant]
\label{lem:L-minorant}
For $t\geq0$, define
\begin{equation}
 L(t):=\frac{t^2}{2(1+2t)},
 \qquad
 L'(t)=\frac{t(1+t)}{(1+2t)^2}.
 \label{eq:L-def}
\end{equation}
Then
\begin{equation}
 0\leq L(t)\leq\mathfrak W(t)
 \qquad(t\geq0).
 \label{eq:L-minorant}
\end{equation}
\end{lemma}

\begin{proof}
For $0\leq t\leq3/4$, one has
$\mathfrak W(t)=t^2/2\geq L(t)$.  For $t\geq3/4$,
\eqref{eq:sawtooth-bounds} gives
\[
 \mathfrak W(t)\geq\frac t4-\frac1{32}.
\]
Moreover,
\[
 \frac t4-\frac1{32}-\frac{t^2}{2(1+2t)}
 =\frac{6t-1}{32(1+2t)}\geq0.
\]
This proves \eqref{eq:L-minorant}.
\end{proof}

Define
\begin{equation}
 \mathcal R_*(\rho,K):=-\int_{(0,\tau)}L(t)\,dg(t).
 \label{eq:R-star-def}
\end{equation}
Because $-dg$ is positive on the decreasing branch,
\begin{equation}
 0\leq\mathcal R_*\leq\mathcal R_{\mathrm{dec}}.
 \label{eq:R-star-order}
\end{equation}
Integration by parts is legitimate at $0$ because
$L(t)g(t)=O(t^{5/3})$.  Hence
\begin{equation}
 \mathcal R_*
 =-L(\tau)g(\tau)+\int_0^\tau L'(t)g(t)\,dt.
 \label{eq:R-star-parts}
\end{equation}
On the outer branch $t=A(z)$, one has
$g(t)\,dt=-2z\,dz$.  Therefore
\begin{equation}
 \mathcal R_*
 =-\frac{\tau^2}{2(1+2\tau)}g(\tau)
 +\int_{\rho K}^{K}
   2z\,\frac{A(z)(1+A(z))}{(1+2A(z))^2}\,dz.
 \label{eq:R-star-action}
\end{equation}
This expression is smooth at every point with $0<\rho<1$ and $K>0$ and has
no floor or wall discontinuity.

\begin{corollary}[Direct retained-remainder bound]
\label{cor:sufficient}
For every $0<\rho<1$ and $K>0$, define
\begin{equation}
 \begin{aligned}
 \mathcal H(\rho,K):={}&
 \frac{\rho^2K^2}{4}
 -\frac{\pi\rho K}{4(1-\rho)}\\
 &+\frac{2\pi\rho K}{\arccos\rho}
   \left(\frac{\tau}{4(1+2\tau)}+\frac3{32}\right)\\
 &+\int_{\rho K}^{K}
   2z\,\frac{A(z)(1+A(z))}{(1+2A(z))^2}\,dz,
 \qquad \tau=KG_1(\rho).
 \end{aligned}
 \label{eq:H-def}
\end{equation}
Then
\begin{equation}
 W(\rho,K)-P(\rho,K)
 \geq \mathcal H(\rho,K)-\mathcal E_{\mathrm{ang}}(\rho,K).
 \label{eq:direct-structural-gap}
\end{equation}
In particular,
$\mathcal E_{\mathrm{ang}}(\rho,K)<\mathcal H(\rho,K)$ implies
$P(\rho,K)<W(\rho,K)$.
\end{corollary}

\begin{proof}
By \eqref{eq:B-def}, \eqref{eq:R-star-action}, and
\[
 \frac\tau4-\frac{\tau^2}{2(1+2\tau)}
 =\frac{\tau}{4(1+2\tau)},
\]
one has $\mathcal H=\mathcal B+\mathcal R_*$.  Now use
\eqref{eq:master-identity},
\eqref{eq:R-star-order}, and $\mathcal R_{\mathrm{osc}}\geq0$.
\end{proof}

\begin{remark}[Handoff]
The companion one-layer estimate proves, whenever at least one shifted
radial layer is active,
\[
 \mathcal E_{\mathrm{ang}}
 <\frac{(1-\rho^2)K^2}{6}-\frac N2,
 \qquad N=\#\{n\geq1:n-1/4<T\}.
\]
Combining this with \eqref{eq:direct-structural-gap} gives one direct smooth
lower bound for $W-P$.  The separate double-sawtooth expansion is not needed
for this implication and is kept outside the live proof as a research
reformulation.
\end{remark}

\end{document}
